Model calibration served as a critical computational validation step. Before the introduction of the Deep Reinforcement Learning (DRL) optimization policies, it was imperative to ensure that the simulation accurately reproduced the empirical ``Status Quo'' of the Municipality of Bacolod at $t=0$. 

\subsection{Addressing Statistical Divergence and Reporting Bias}
A significant algorithmic challenge identified during the initialization phase was the severe statistical divergence between the compliance rates self-reported by local officials and the macro-level reality established by the Municipal Environment and Natural Resources Office (MENRO) audit. Aggregating the micro-level interview data suggested a municipal compliance rate exceeding 50\%—with localities like Binuni and Demologan reporting unrealistic peaks of 95\% and 60\%, respectively. However, the objective municipal-wide waste characterization study conducted by MENRO confirmed a functional segregation rate of only 10\% to 15\% (refer to Appendix \ref{appendix:menro_interview}).

Utilizing the unverified, inflated self-reports would introduce severe model discrepancy and render the AI's optimization tasks computationally trivial. This divergence quantitatively captured the ``Act vs.\ Reality'' gap described by \cite{Yazawa2025} and the phenomenon of Socially Desirable Responding (SDR) \citep{Geiger2021}, where respondents conflate theoretical policy awareness with actual on-the-ground practice. Consequently, the micro-level survey data was deemed methodologically unreliable for establishing the initial state space ($S_0$) of the simulation, and the MENRO audit ($\approx 12.5\%$) was selected as the ground-truth calibration target.

\subsection{Algorithmic Correction of Initial Compliance States}
To ensure the Agent-Based Model reflected true baseline conditions, the initialization parameters in \texttt{barangay\_config.py} were systematically calibrated downward. Table \ref{tab:calibration_data} details the algorithmic correction applied to the raw interview data to align the local nodes with the municipal reality.

\begin{table}[htbp]
    \centering
    \caption{Algorithmic Calibration of Self-Reported vs.\ Baseline Compliance Rates}
    \label{tab:calibration_data}
    \small
    \begin{tabular}{l c c p{6.5cm}}
        \toprule
        \textbf{Barangay Node} & \textbf{Self-Reported} & \textbf{Calibrated ($t=0$)} & \textbf{Primary Calibration Justification} \\
        \midrule
        Brgy Babalaya     & 100\% & 15.0\% & Adjusted for lack of persistent funding/monitoring. \\
        Brgy Binuni       & 95\%  & 15.0\% & Decoupled theoretical awareness from physical practice. \\
        Brgy Mati         & 70\%  & 14.0\% & Corrected for high habit decay rate in agrarian zones. \\
        Brgy Liangan East & 65\%  & 14.0\% & Smoothed transient project-based spikes (Eco-bricks). \\
        Brgy Demologan    & 60\%  & 14.0\% & Standardized to the municipal empirical baseline. \\
        Brgy Ezperanza    & 20\%  & 13.5\% & Aligned with the middle-class control group variance. \\
        Brgy Poblacion    & 2\%   & 2.0\%  & \textbf{Anchor Point} (Aligned perfectly with MENRO Audit). \\
        \bottomrule
    \end{tabular}
    
    \vspace{1ex}
    \raggedright
    \footnotesize
    \textit{Note: ``Self-Reported'' refers to raw rates claimed during Key Informant Interviews (Appendix \ref{ch:appendix_b}). ``Calibrated'' refers to the initialized compliance state within the ABM, mathematically tuned to match the aggregate municipal segregation rate of $\approx 12.5\%$ verified by MENRO.}
\end{table}

This calibration was executed through three specific analytical pivots:
\begin{itemize}
    \item \textbf{Decoupling Awareness from Practice (Binuni):} Officials reported a near-perfect compliance rate of 95\%. However, initializing the simulation at this level resulted in immediate global compliance exceeding 40\%, bypassing the behavioral friction the DRL was designed to solve. The simulation logic posited that the 95\% reflected \textit{awareness}. As noted by \cite{Paigalan2025}, positive attitudes frequently coexist with poor practices, creating an ``Intention-Action Gap'' that the model explicitly corrected by dropping the initialization to 15\%.
    \item \textbf{Smoothing Transient Spikes (Liangan East):} The reported 65\% compliance was driven by a temporary Eco-brick initiative. Because project-based compliance is historically transient without continuous fiscal injection \citep{Camarillo2021}, the system initialized Liangan East at 14\% to model the habituated ``Status Quo'' before the DRL agent assumed budgetary control.
    \item \textbf{The Urban Anchor Point (Poblacion):} The only data point where qualitative self-reporting matched the quantitative MENRO audit was Barangay Poblacion (2\%). Because both metrics converged, this served as the simulation's \textit{Anchor Point}. This validated the ``Urban Resource Trap'' theory: high population density and anonymity dissolve social pressure, defaulting to non-compliance regardless of existing legal frameworks \citep{Villanueva2021}.
\end{itemize}

\subsection{Stochastic Seeding and Mathematical Validation}
To reconcile the remaining micro-level bias with the verified macro-level ground truth, the initialized states for the remaining barangays were assigned stochastically. Rather than falsely assuming homogeneous behavior by hard-coding a static 12.5\% across all 4,795 agents, the initialization parameters were drawn from a bounded uniform distribution, $U(0.11, 0.15)$. 

In Agent-Based Modeling, utilizing stochastic initialization bounded by verified empirical data is a standard practice for uncertainty quantification \citep{McCulloch2022, Fischer2025}. The mathematical expected value ($\mu$) of this specific distribution anchors tightly to the $\approx 12.5\%$ MENRO baseline. This technique allowed the model to introduce realistic heterogeneous variance (noise) among the different barangay grids while strictly preserving the global starting state.

This distributional parameterization was critical for the structural validity of the Deep Reinforcement Learning environment. As \cite{Ding2024} emphasize, ABMs are highly sensitive to initial parameterization; failing to align these vectors with empirical reality prevents the simulation from yielding physically meaningful quantities. By anchoring the randomization to this conservative baseline, headless simulation benchmarks successfully stabilized at a Global Average Compliance of $\approx 12.5\%$ at $t = 0$. This mathematically proved the structural validity of the ABM, ensuring the subsequent DRL agent was optimizing its policies against a rigorously realistic scenario rather than a mathematically trivial one.